\subsection{The polygon associated to a regular region}
We define a polygonal region of $\mathbb{R}^2$ as a space locally modeled on \say{wedges}. In particular, let $\arg \colon \mathbb{R}^2 \to [0, 2\pi)$ be a branch of the argument function, and for $\theta \in (-\pi, \pi]$, define the \emph{closed sector with exterior angle $\theta$} to be the subset
\begin{equation}
  W_{\theta} = \{z \mid 0 \leq \arg(z) \leq \pi + \theta \} \subseteq \mathbb{R}^2
\end{equation}
For a subset $S \subseteq \mathbb{R}^2$ and $p \in S$, define a \emph{sector chart} at $p$ to be a tuple $(U, V, \phi)$ consisting of an open neighborhood $U$ of $p$, an open neighborhood $V$ of the origin, and an isometry
\begin{equation}
  \phi \colon U \cap S \to V \cap W_\theta
\end{equation} 
or some $\theta \in (-\pi, \pi]$ which takes $p$ to the origin.
If such a chart exists, then $\theta = \theta_S(p)$ is uniquely determined, and is said to be the \emph{exterior angle} at $p$. With this, we define a \emph{polygonal region} to be a subset $P \subseteq \mathbb{R}^2$ which admits a sector chart at each point $p \in P$. The exterior angle function partitions the region into three sets:
\begin{enumerate}
  \item a set $\operatorname{int}(P)$ of  \emph{interior points}, whose exterior angles are all equal to $\pi$,
  \item a set $\operatorname{bd}(P)$ of \emph{boundary points}, whose exterior angles are all equal to $0$, and
  \item a set $\operatorname{cr}(P)$ of \emph{corner points} or \emph{vertices}, whose exterior angles are not equal to $\pi$ or $0$.
\end{enumerate}
Note that the interior and boundary of a polygonal region as defined above agree with the usual topological definitions. Within a sector chart, there must be at most one corner point, so $\operatorname{cr}(P)$ is isolated. Thus, for a bounded polygonal region, there exist finitely many corners. We reformulate a well-known special case of the Gauss-Bonnet theorem in terms of the exterior angle function.
\begin{theorem}\label{thm:gauss-bonnet}
  Let $P \subseteq \mathbb{R}^2$ be a connected, bounded, polygonal region with $b$ boundary components. Then
  \begin{equation}
    \sum_{p \, \in \, \operatorname{cr}(P)} \!\!\! \theta_P(p) = 2(2-b)\pi. 
  \end{equation}
\end{theorem}
\begin{proof}
  See \cite{docarmo1976}.
\end{proof}